% !PRINTED 40
\weberpage{40}
%<<CONT

\[
(x-\alpha_1)(x-\alpha_2)=x^2-(\alpha_1+\alpha_2)x+\alpha_1\alpha_2,
\]
\[
\begin{aligned}
 (x-\alpha_1)(x-\alpha_2)(x-\alpha_3)&=x^3-(\alpha_1+\alpha_2+\alpha_3)x^2\\
&\quad +(\alpha_1\alpha_2+\alpha_1\alpha_3+\alpha_2\alpha_3)x-\alpha_1\alpha_2\alpha_3.
\end{aligned}
\]
Nehmen wir die Richtigkeit unseres Bildungsgesetzes bei $n-1$ Factoren als bewiesen an und setzen
\[
 (x-\alpha_1)(x-\alpha_2)\cdots(x-\alpha_{n-1})=x^{n-1}+a'_1x^{n-2}+\cdots+a'_{n-1},
\]
so findet man durch Multiplication mit $x-\alpha_n$:
\begin{equation}
\begin{aligned}
 a_1&=a'_1-\alpha_n,\\
 a_2&=a'_2-\alpha_n a'_1,\\
 a_3&=a'_3-\alpha_n a'_2,\\
&\cdots\\
 a_\nu&=a'_\nu-\alpha_n a'_{\nu-1},\\
&\cdots\\
 a_n&=-\alpha_n a'_{n-1},
\end{aligned}
\tag{3}
\end{equation}
und daraus ist die Richtigkeit der Formeln (2) unmittelbar ersichtlich.

Es ist von Wichtigkeit, die Anzahl der Glieder zu bestimmen, die in jeder der Summen (2) vorkommen. Wir bezeichnen die Anzahl der Terme, die in der Summe $(-1)^\nu a_\nu$ vorkommen, mit $B_\nu^{(n)}$; es ist die Anzahl der Combinationen ohne Wiederholung von $n$ Elementen zur $\nu$ten Classe (d. h. von je $\nu$ verschiedenen der $n$ Elemente). Um sie zu bestimmen, denke man sich zunächst die $B_{\nu-1}^{(n)}$ Combinationen zur $(\nu-1)$ten Classe gebildet. Aus jeder dieser Combinationen kann man durch Hinzufügung je eines der fehlenden $n-\nu+1$ Elemente $n-\nu+1$ Combinationen zur $\nu$ten Classe ableiten. Auf diese Art aber wird jede Combination $\nu$mal, nämlich durch Hinzufügung jedes ihrer Elemente gebildet, so dass man die Relation
\begin{equation}
 B_\nu^{(n)}=B_{\nu-1}^{(n)}{n-\nu+1\over \nu}
\tag{4}
\end{equation}
erhält, während $B_1^{(n)}$ offenbar den Werth $n$ hat. Wenn man also die aus (4) folgenden Gleichungen
\begin{equation}
\begin{aligned}
 B_1^{(n)}&=n,\\
 B_2^{(n)}&=B_1^{(n)}{n-1\over 2},\\
&\cdots\\
 B_\nu^{(n)}&=B_{\nu-1}^{(n)}{n-\nu+1\over \nu}
\end{aligned}
\tag{5}
\end{equation}
multiplicirt, so folgt:
%CONT>>
